\section{Introduction}
\label{submission:intro}

The rapid expansion of pre-trained foundation models has led to a major paradigm shift in machine learning: rather than training massive networks from scratch, practitioners fine-tune pre-trained models on specific downstream tasks \cite{dosovitskiy2020vit, wortsman2022model}. While highly effective, this process yields a large collection of disjoint, task-specific expert models. Managing, storing, and serving a separate copy of a large model for every task is computationally prohibitive. Consequently, the field of \emph{weight-space model merging} has emerged as an attractive solution \cite{ilharco2023editing, wortsman2022model}. By mathematically interpolating the parameter weights of task-specific experts into a single shared backbone, model merging consolidates multi-task capabilities without requiring any additional training.

However, the primary bottleneck in model merging is \emph{spatial weight-space interference} \cite{yadav2023ties}. When experts are fine-tuned on highly orthogonal or conflicting domains, they learn parameter updates that move in opposing directions or occupy distinct regions of representation. Standard linear interpolation (Task Arithmetic) \cite{ilharco2023editing} averages these parameter updates, which inevitably dilutes task-specific features and results in catastrophic representational collapse.

In response to this interference, the deep learning community has proposed various advanced methods. Recent literature has introduced test-time adaptation (TTA) frameworks that optimize hundreds of layer-wise scaling coefficients and magnitude-pruning boundaries concurrently on calibration data \cite{yu2024dare, yadav2023ties}. These methods often require backpropagation loops, auxiliary loss networks, or continuous parameter search spaces. 

\begin{figure}[t]
  \centering
  \includegraphics[width=0.45\textwidth]{../qmerge_vs_baselines.png}
  \caption{Conceptual overview of Exclusive Parameter Merging (EPM). Standard average-based methods (Task Arithmetic, TIES-Merging) blend updates across different experts, creating massive weight conflicts and representation collapse. EPM mitigates this at the coordinate level: each weight parameter is primarily assigned to the expert with the dominant scaled update magnitude while retaining a background fraction of the standard blend ($\gamma=0.2$), suppressing spatial weight-space interference while preserving network coherence across layers.}
  \label{fig:epm_concept}
\end{figure}

From a practical perspective, we argue that this added parameterization can often be simplified. Highly parameterized test-time optimization can suffer from the \textbf{Overfitting-Optimizer Paradox}: continuous high-dimensional tuning on small calibration or validation datasets minimizes the proxy optimization loss (such as Shannon entropy) but can overfit transductively, destroying generalizable multi-task features. This results in fragile models that perform poorly under slight task or domain shifts.

This leads us to a central inquiry: \emph{How can we mitigate weight-space interference directly during parameter composition, thereby reducing the dependency on high-dimensional post-hoc optimizations?}

To answer this, we propose \textbf{Exclusive Parameter Merging (EPM)}, a conceptually straightforward, training-free routing protocol that mitigates spatial weight-space interference. While traditional methods perform uniform average weight blending, EPM implements \textbf{Soft Exclusive Parameter Allocation (Soft-EPA)} as a spatially selective regularizer: for every parameter coordinate across the model, EPM identifies the dominant task expert based on scaled absolute magnitude. It routes the dominant update at full strength, while attenuating non-dominant updates by a coherence retention factor $\gamma = 0.2$. This balanced approach leverages the benefits of coordinate exclusivity to shield complex representation channels while allowing a small background linear blend to preserve network-wide activation manifold alignment. We explicitly acknowledge that while standard average-based weight sharing can achieve slightly higher joint average performance under high sparsity settings (due to the natural spatial decoupling induced by pruning), Soft-EPA serves as an essential spatial safeguard that prevents individual expert collapse and raises the worst-case multi-task performance floor.

To optimize the scaling factors while avoiding the overfitting-optimizer paradox, we introduce \textbf{Task-Level Coefficient Tuning (TLC-Tune)}. Instead of hundreds of layer-wise variables, TLC-Tune restricts the search space to only $K$ global scaling factors (one per task expert). Since this $K$-dimensional search space is extremely low-dimensional and stable, we optimize it using a simple zero-order (1+1) Evolution Strategy on a stable offline validation split of 128 samples per task. TLC-Tune is perfectly stable, training-free, and highly robust to transductive noise.

We conduct exhaustive multi-task evaluations on a Vision Transformer (ViT-Tiny) backbone across four highly conflicting domains: MNIST, FashionMNIST, CIFAR-10, and SVHN. Our empirical findings demonstrate that:
\begin{itemize}
  \item \textbf{Mitigation of Representational Collapse:} Under severe domain conflicts, standard average-based methods collapse to near-random accuracies. EPM successfully mitigates representational collapse. By incorporating Task Vector Standardization, EPM resolves the "Rich Task" Dominance Trap, leveling the playing field and elevating simple task (MNIST, FashionMNIST) accuracies from random guessing to robust, highly functional states.
  \item \textbf{Empirical Outperformance:} EPM consistently outclasses traditional linear averaging, TIES-Merging, and Random Tensor Routing across target sparsities of 0\% and 50\%, while achieving performance highly competitive with the state-of-the-art DARE baseline, reaching up to \textbf{46.19\%} joint mean accuracy under dense merging with extremely balanced multi-task validation.
  \item \textbf{Deconstructing Overfitting vs. Optimization Failure:} We evaluate EPM against highly parameterized layer-group-wise optimization baselines (AdaMerging and ZipMerge). Through a systematic 500-step optimization study, we empirically demonstrate that the failure of these high-dimensional methods on small calibration datasets is a dual catastrophe: they suffer from absolute optimization failure (under-convergence) under zero-order search, while remaining highly vulnerable to transductive overfitting if continuously tuned. In contrast, TLC-Tune's minimalist 4-dimensional global coefficient space is easily optimized and immune to transductive noise, ensuring flawless test set generalization.
\end{itemize}
Our core routing implementation requires fewer than 10 lines of PyTorch, illustrating that direct, mathematically sound coordinate-level operators can offer a highly competitive and robust alternative to highly parameterized optimization pipelines.
